\chapter{Contexte et Problématique}

\section{Introduction}

Ce premier chapitre pose les bases du projet LudoHeritage. Il présente le contexte culturel et
numérique dans lequel s'inscrit le projet, identifie la problématique centrale à laquelle il
répond, analyse les solutions existantes, et définit clairement les objectifs visés.

% ─────────────────────────────────────────────────────────────────────────────
\section{Les jeux de société traditionnels : un patrimoine en danger}
% ─────────────────────────────────────────────────────────────────────────────

Les jeux de société traditionnels représentent une forme unique de patrimoine culturel immatériel.
Contrairement aux monuments ou aux œuvres d'art, ce patrimoine \og{}vivant\fg{} se transmet de
génération en génération, portant avec lui les valeurs, les modes de pensée et l'histoire des
peuples qui les ont créés.

L'UNESCO, dans sa Convention pour la Sauvegarde du Patrimoine Culturel Immatériel (2003), reconnaît
explicitement les \textit{pratiques sociales, rituels et événements festifs} -- dont font partie les
jeux traditionnels -- comme des composantes essentielles du patrimoine mondial à préserver.

Parmi les jeux recensés dans LudoHeritage, on trouve des exemples qui illustrent parfaitement cette
richesse culturelle~:

\begin{itemize}
    \item \textbf{L'Awale} (Afrique de l'Ouest, \textasciitilde 3000 av. J.-C.) est bien plus
          qu'un jeu de stratégie. Il est un outil d'enseignement du calcul mental, un rituel de
          socialisation communautaire, et un symbole de l'ingéniosité africaine.

    \item \textbf{Le Toguz Korgool} (Kirghizstan), inscrit au patrimoine immatériel de l'UNESCO,
          est une institution nationale jouée lors des fêtes et transmise par les aînés aux jeunes
          générations.

    \item \textbf{Le Fanorona} (Madagascar) est indissociable de l'identité malgache. La légende
          raconte que le roi Ralambo y jouait au XVII\ieme{} siècle, et il est aujourd'hui classé
          patrimoine culturel national.

    \item \textbf{L'Alquerque} (Monde arabe, XI\ieme{} siècle) a été documenté dans le
          \textit{Libro de los Juegos} du roi Alphonse X de Castille. Il est le précurseur direct
          des Dames modernes.
\end{itemize}

Malgré cette valeur indéniable, ces jeux souffrent d'un effacement progressif. La mondialisation
des loisirs, la domination des jeux vidéo et l'absence de documentation numérique accessible
contribuent à leur disparition progressive dans la mémoire collective.

% ─────────────────────────────────────────────────────────────────────────────
\section{La problématique}
% ─────────────────────────────────────────────────────────────────────────────

L'analyse du paysage numérique actuel révèle plusieurs lacunes importantes~:

\subsection{Absence d'une plateforme francophone dédiée}

Les rares ressources numériques existantes sont quasi-exclusivement en anglais. Un public
francophone -- notamment au Maroc, en Afrique subsaharienne ou au Maghreb -- n'a pas accès à
une plateforme structurée, interactive et dans sa langue pour explorer ces jeux.

\subsection{Dispersion de l'information}

Les informations sur les jeux traditionnels sont éparpillées entre des articles académiques,
des encyclopédies en ligne, des forums spécialisés et des sites culturels locaux. Il n'existe
pas de source unique, centralisée et fiable qui regroupe règles, histoire et contexte culturel.

\subsection{Manque d'outils pédagogiques interactifs}

Les enseignants, les étudiants et les curieux qui souhaitent découvrir ces jeux n'ont pas accès
à des outils interactifs~: pas de cartes pour situer géographiquement les jeux, pas de frises
chronologiques, pas d'assistants pour répondre à leurs questions, pas de quiz pour tester leurs
connaissances.

\subsection{Isolement des communautés}

Les passionnés de jeux traditionnels n'ont pas d'espace numérique communautaire en français pour
partager leurs découvertes, poser des questions et enrichir collectivement la connaissance.

\begin{figure}[H]
    \centering
    \begin{tabular}{|p{5cm}|p{8cm}|}
        \hline
        \textbf{Problème identifié} & \textbf{Impact} \\
        \hline
        Aucune plateforme francophone & Public de 300+ millions de locuteurs mal desservi \\
        \hline
        Information dispersée & Temps de recherche élevé, risque d'informations inexactes \\
        \hline
        Pas d'outils interactifs & Apprentissage passif, engagement faible \\
        \hline
        Pas de communauté en ligne & Connaissances non partagées, perte de savoirs \\
        \hline
    \end{tabular}
    \caption{Synthèse des problèmes identifiés}
    \label{tab:problemes}
\end{figure}

% ─────────────────────────────────────────────────────────────────────────────
\section{État de l'art}
% ─────────────────────────────────────────────────────────────────────────────

Avant de concevoir LudoHeritage, une analyse des solutions existantes a été effectuée~:

\begin{description}
    \item[\textbf{Ludii (ludii.games)}] Plateforme de recherche académique développée par
    l'Université de Maastricht dans le cadre du \textit{Digital Ludeme Project} (financement ERC
    de 2 millions d'euros). Elle propose un système de description formelle des jeux (ludèmes) et
    une base de données de plusieurs milliers de jeux. Cependant, elle est \textbf{entièrement en
    anglais}, très technique, destinée principalement aux chercheurs, et ne propose pas d'assistant
    IA ni de forum communautaire.

    \item[\textbf{BoardGameGeek}] La plus grande base de données en ligne de jeux de société.
    Elle est riche en contenu mais concentrée sur les jeux modernes, \textbf{en anglais uniquement},
    et ne propose aucun outil d'exploration culturelle ni d'assistant.

    \item[\textbf{Wikipedia / Articles culturels}] Contenu riche mais non structuré, sans
    interactivité, sans visualisation géographique ni chronologique.
\end{description}

Ce panorama confirme qu'il existe un réel \textbf{vide numérique} pour une plateforme francophone,
interactive et culturellement engagée dédiée aux jeux traditionnels.

% ─────────────────────────────────────────────────────────────────────────────
\section{Objectifs du projet}
% ─────────────────────────────────────────────────────────────────────────────

LudoHeritage se fixe les objectifs suivants~:

\begin{enumerate}
    \item \textbf{Centraliser} un catalogue de vingt jeux traditionnels soigneusement sélectionnés,
          avec leurs règles, leur histoire, leurs mécaniques et leur valeur culturelle.

    \item \textbf{Visualiser} le patrimoine ludique à travers une carte mondiale interactive et
          une frise chronologique allant de l'Antiquité à nos jours.

    \item \textbf{Assister} les utilisateurs grâce à un chatbot IA (LudoBot) capable de répondre
          à des questions sur les jeux, de faire des recommandations et d'expliquer des règles.

    \item \textbf{Éduquer} via un module de quiz interactif avec classement, un outil de
          comparaison entre deux jeux, et des tutoriels pas-à-pas pour chaque jeu.

    \item \textbf{Engager} la communauté à travers un forum de discussion, un système de favoris
          et un profil utilisateur personnalisé.

    \item \textbf{Respecter} les valeurs culturelles et religieuses en écartant tout contenu
          associé à des pratiques contraires à l'éthique islamique (jeux de hasard rituels,
          contenu à caractère religieux polythéiste).
\end{enumerate}

% ─────────────────────────────────────────────────────────────────────────────
\section{Conclusion}
% ─────────────────────────────────────────────────────────────────────────────

Ce chapitre a mis en lumière la nécessité d'un projet comme LudoHeritage. Face à la disparition
progressive du patrimoine ludique mondial et à l'absence de solutions numériques francophones
adaptées, ce projet propose une réponse concrète et complète. Le prochain chapitre détaillera
comment ce projet a été analysé, conçu et structuré pour atteindre ces objectifs.
